% Owner: Sumaiya. Rubric: Problem Description ("well described, clear motivation, what it aims to solve").
\section{Introduction}\label{sec:intro}

Wikipedia is edited around the clock. Every edit on every Wikimedia wiki appears within seconds in a
public live feed~\cite{wikimedia2026eventstreams}, and after keeping only edits and page creations that is
about 20 events per second, or roughly 1.7 million a day. Some come from registered volunteers, some from
logged-out editors, and many from bots: automated accounts that fix links, import data or revert
vandalism. Volunteer \emph{patrollers} review recent changes for vandalism and
mistakes~\cite{wikipedia2026patrol}. Their time is limited, and the stream does not say where it is best
spent.

We ask one question: \emph{where does editing pressure concentrate across Wikipedia's language editions,
how much of it comes from bots, and are spikes in reader attention followed by edit surges?} By editing
pressure we mean the number of edits per wiki over time that may need review. The question has three
parts because patrollers decide on three levels: which wikis need the most attention, how much of the
volume is flagged bot work that is hidden from their default review list, and whether a sudden rise in readers is an early sign of
a rush of edits.

The audience is patrollers and wiki administrators deciding where to focus review effort. Each wiki's
community grants the bot flag only to approved bots, and flagged edits are hidden from the default list of
recent changes, so the edits left for human review are mostly the unflagged ones. Reader attention matters
because an article in the news gets many readers at once, and its edits may rise at the same time.

To answer the question we built an end-to-end pipeline that runs on one laptop with one command:
\begin{itemize}
  \item a \textbf{stream path} that reads every edit live through Kafka, validates it, removes user names
    and splits bots from humans;
  \item a \textbf{batch path} that loads the hourly pageview files with Airflow~\cite{wikimedia2026pageviews};
  \item a \textbf{data lake} in MinIO and a \textbf{star schema} in DuckDB, built and tested with dbt;
  \item a \textbf{Metabase dashboard} that answers each part of the question over a 15-day window
    (28 Oct -- 11 Nov 2026, UTC).
\end{itemize}
